\chapter{SOLUCIÓN PROPUESTA}
\label{ch:solucion}

Este capítulo describe, de forma general, la solución diseñada antes de entrar
al detalle teórico y de implementación. Se propone un robot móvil diferencial
de bajo costo que integra percepción láser 2D, odometría por encoders y control
PID embebido, sobre una arquitectura distribuida en dos entornos; robot y
estación de trabajo, comunicados mediante \gls{ROS2}. El objetivo es que el
sistema sea capaz de mapear, navegar y explorar entornos interiores de forma
autónoma.

\utemFigurePropia{2.Texto/Imag/arquitectura_alto_nivel.png}{1.0\textwidth}
{Diagrama de bloques de la solución propuesta.}
{fig:arquitectura_alto_nivel}

La Fig.~\ref{fig:arquitectura_alto_nivel} organiza la solución en cuatro
dominios funcionales, identificados por color:

\begin{description}
\item[Sensado.] Los \emph{encoders} de cuadratura miden el giro de cada rueda
y el \gls{LIDAR} 2D percibe el entorno; son las dos fuentes de información del
sistema.
\item[Control de bajo nivel.] El \gls{Arduino} lee los encoders y ejecuta el
control PID de velocidad de cada rueda, accionando los motores DC de tracción
diferencial.
\item[Sistema computacional del robot.] La Raspberry Pi~4 integra \gls{ROS2},
estima la odometría a partir de las velocidades de rueda y actúa de puente
entre el hardware embebido y la estación de trabajo.
\item[Mapeo y navegación.] En la estación de trabajo, \texttt{slam\_toolbox}
construye el mapa y estima la pose, Nav2 planifica y controla la navegación,
el explorador de fronteras genera metas autónomas y RViz2 permite la
supervisión.
\end{description}

El flujo de datos es el siguiente: los encoders realimentan al Arduino, que
cierra el lazo PID y entrega los conteos acumulados; la interfaz
\texttt{DiffDriveArduino} convierte esos conteos en estados de rueda y
\texttt{diff\_drive\_controller} publica \texttt{/odom} y la transformación
\texttt{odom}$\rightarrow$\texttt{base\_link}. La Raspberry Pi publica además
los barridos \texttt{/scan}; en la estación de trabajo,
\texttt{slam\_toolbox} consume \texttt{/scan} y obtiene la pose odométrica
mediante TF, en lugar de suscribirse directamente a \texttt{/odom}, para
estimar la trayectoria y construir el mapa, que posteriormente alimenta a
Nav2 y al explorador de fronteras.